\documentclass[10pt,twoside]{article}\usepackage[]{graphicx}\usepackage[dvipsnames,svgnames,table]{xcolor}
% maxwidth is the original width if it is less than linewidth
% otherwise use linewidth (to make sure the graphics do not exceed the margin)
\makeatletter
\def\maxwidth{ %
  \ifdim\Gin@nat@width>\linewidth
    \linewidth
  \else
    \Gin@nat@width
  \fi
}
\makeatother

\definecolor{fgcolor}{rgb}{0.345, 0.345, 0.345}
\newcommand{\hlnum}[1]{\textcolor[rgb]{0.686,0.059,0.569}{#1}}%
\newcommand{\hlsng}[1]{\textcolor[rgb]{0.192,0.494,0.8}{#1}}%
\newcommand{\hlcom}[1]{\textcolor[rgb]{0.678,0.584,0.686}{\textit{#1}}}%
\newcommand{\hlopt}[1]{\textcolor[rgb]{0,0,0}{#1}}%
\newcommand{\hldef}[1]{\textcolor[rgb]{0.345,0.345,0.345}{#1}}%
\newcommand{\hlkwa}[1]{\textcolor[rgb]{0.161,0.373,0.58}{\textbf{#1}}}%
\newcommand{\hlkwb}[1]{\textcolor[rgb]{0.69,0.353,0.396}{#1}}%
\newcommand{\hlkwc}[1]{\textcolor[rgb]{0.333,0.667,0.333}{#1}}%
\newcommand{\hlkwd}[1]{\textcolor[rgb]{0.737,0.353,0.396}{\textbf{#1}}}%
\let\hlipl\hlkwb

\usepackage{framed}
\makeatletter
\newenvironment{kframe}{%
 \def\at@end@of@kframe{}%
 \ifinner\ifhmode%
  \def\at@end@of@kframe{\end{minipage}}%
  \begin{minipage}{\columnwidth}%
 \fi\fi%
 \def\FrameCommand##1{\hskip\@totalleftmargin \hskip-\fboxsep
 \colorbox{shadecolor}{##1}\hskip-\fboxsep
     % There is no \\@totalrightmargin, so:
     \hskip-\linewidth \hskip-\@totalleftmargin \hskip\columnwidth}%
 \MakeFramed {\advance\hsize-\width
   \@totalleftmargin\z@ \linewidth\hsize
   \@setminipage}}%
 {\par\unskip\endMakeFramed%
 \at@end@of@kframe}
\makeatother

\definecolor{shadecolor}{rgb}{.97, .97, .97}
\definecolor{messagecolor}{rgb}{0, 0, 0}
\definecolor{warningcolor}{rgb}{1, 0, 1}
\definecolor{errorcolor}{rgb}{1, 0, 0}
\newenvironment{knitrout}{}{} % an empty environment to be redefined in TeX

\usepackage{alltt}
\usepackage[marginparsep=1em]{geometry}
\geometry{lmargin=1.0in,rmargin=1.0in, bmargin=1.2in,  tmargin=1.2in}
\usepackage[dvipsnames,svgnames,table]{xcolor}
\usepackage{graphicx}
\usepackage{amssymb}
\usepackage{epstopdf}
\usepackage{verbatim}
\usepackage{enumerate}
\usepackage{bm}
\usepackage{amsthm}
\usepackage{float}
\usepackage{amsmath}
\usepackage{fancyhdr}
\usepackage{hyperref}
\usepackage{mathtools}

%%%%  SHORTCUT COMMANDS  %%%%
\newcommand{\ds}{\displaystyle}
\newcommand{\Z}{\mathbb{Z}}
\newcommand{\T}{\mathcal{T}}
\newcommand{\arc}{\rightarrow}
\newcommand{\R}{\mathbb{R}}
\newcommand{\RP}{\mathbb{R}(+)}
\newcommand{\Rs}{\mathbb{R}^{**}}
\newcommand{\C}{\mathbb{C}}
\newcommand{\E}{\mathbb{E}}
\newcommand{\B}{\mathcal{B}}
\newcommand{\PX}{\mathcal{P}(X)}
\newcommand*\rot{\rotatebox{90}}
\newcommand*\OK{\ding{51}}
\newcommand{\N}{\mathbb{N}}
\newcommand{\Q}{\mathbb{Q}}
\newcommand{\Answer}{\vspace{2mm}\textbf{\underline{Answer}}\\}
\newcolumntype{L}{>{$}l<{$}}
\newcolumntype{C}{>{$}c<{$}}
\newcommand{\GL}{\text{GL}}
\newcommand{\SL}{\text{SL}}
\newcommand{\Var}{\text{Var}}

\newcommand{\stirling}[2]{\genfrac{\{}{\}}{0pt}{}{#1}{#2}}

\pagestyle{fancy}
\fancyhf{}
\renewcommand{\sectionmark}[1]{\markright{\thesection.\ #1}}
\lhead{\fancyplain{}{}}
\fancyhead[RE,RO]{Math 4451, Spring 2026}
\fancyfoot[RE,RO]{\thepage}
\IfFileExists{upquote.sty}{\usepackage{upquote}}{}
\begin{document}
\begin{flushright}
\begin{minipage}{.33\textwidth}
\rightline{YOUR NAME}
\rightline{\href{mailto:YOUR EMAIL}{YOUR EMAIL}}
\rightline{\today}
\end{minipage}
\end{flushright}

\begin{center}
{\large{\textbf{Homework 7}}}
\end{center}

\begin{enumerate}
  \item You are testing the lifespan of a new type of computer part. Let $X_1, X_2, \ldots, X_n$ be independent and identically distributed lifespans of a sample of $n$ of these parts. A common model for the lifespan of an electrical component is the exponential distribution. We will parameterize this distribution by its unknown mean, $\theta$. The probability density function (PDF) is:
  $$f_X(x; \theta) = \frac{1}{\theta} e^{-x/\theta}, \quad x > 0$$
  \begin{enumerate}
    \item (1 point) \label{prob:expMin} Suppose an impatient technician decides to estimate the mean lifespan by only recording the lifespan of the very first computer part to fail. That is, let $Y = \min(X_1, X_2, \ldots, X_n)$, and we want to find an unbiased estimator of $\theta$ based only on $Y$ (you may use the densities for order statistics derived last semester).
    \item (1 point) Calculate the variance of your estimator in part~\ref{prob:expMin}.
    Note that it may be helpful to compare the density of your estimate and compare to known distributions.
    % \item Find an unbiased estimator of $\lambda$ based only on $Y = \min\{X_1, \ldots, X_n\}$. (Hint, consider the order statistics formulas from last semester).
    \item (1 point) Now derive the maximum likelihood estimator $\hat{\theta}$ for $\theta$.
    \item (1 point) \label{prob:expBias} Calculate the Bias of the MLE.
    \item (1 point) Calculate the Variance of the MLE.
    \item (1 point) Is the MLE a consistent estimator of $\theta$? Why or why not.
    \item (1 point) Which estimator do you prefer and why?
    \item (1 point) Let $\tilde{\theta}$ be some other estimator with the property that $E[\tilde{\theta}] = E[\hat{\theta}]$ (i.e., another estimator with the same bias from part~\ref{prob:expBias}). What is the absolute lowest possible variance for the estimate $\tilde{\theta}$? Does $\hat{\theta}$ achieve this lower bound?
  \end{enumerate}
  \item Suppose you are helping an engineer analyze the magnitude of background static noise in a new wireless communication channel. You measure the noise amplitude over $n$ independent trials, resulting in the sample $X_1, X_2, \ldots, X_n$.

  You model the noise amplitude using a Rayleigh distribution, which depends on the parameter $\theta$, and has pdf:
  $$
  f(x; \theta) = \frac{x}{\theta} \exp\Big\{-\frac{x^2}{2\theta}\Big\}, \quad x > 0.
  $$
  For your own reference, note that $E[X^2] = 2\theta$, and $\Var(X^2) = 4\theta^2$.
  \begin{enumerate}
    \item (1 point) Write down the log-likelihood function, $\ell(\theta)$, for all $n$ observations.
    \item (1 point) Find the MLE of $\theta$ for this model.
    \item (1 point) Calculate the bias and variance of the MLE. Is the MLE a consistent estimator?
    \item (1 point) Calculate the Fisher Information $\mathcal{I}_n(\theta)$ for the entire sample (Note: you can use the expected value of the second derivative of the log-likelihood).
    \item (1 point) Let $\tilde{\theta}$ be some other estimator with the property that $E[\tilde{\theta}] = E[\hat{\theta}]$ (i.e., another estimator with the same bias from part (c)). What is the absolute lowest possible variance for the estimate $\tilde{\theta}$? Does $\hat{\theta}$ achieve this lower bound?
  \end{enumerate}
\end{enumerate}


\end{document}
